% Table S13: Evidence scope and next validation
\begin{table}[htbp]
\centering
\scriptsize
\caption{Evidence scope and next validation. The computational results support prioritization and structural hypotheses within the stated cohorts; they do not substitute for biochemical, biophysical, phenotypic, or clinical measurements. The table separates the primary computational triage estimand from secondary diagnostics, technical checks, and unobserved biological quantities.}
\label{tab:s13_evidence_scope}
\begin{tabularx}{\textwidth}{>{\raggedright\arraybackslash}p{0.19\textwidth} >{\raggedright\arraybackslash}p{0.27\textwidth} >{\raggedright\arraybackslash}X}
\toprule
Evidence layer & Supported interpretation & Next validation needed \\
\midrule
Docking-RRS & Relative retention of empirical docking scores across the selected mutant panel; primary computational triage only & Matched biochemical activity measurements for wild-type and mutant PfDHFR/PfCRT \\
ACSI/PNS & Cohort-relative chemical-space and network-weighted ranking dimensions & Larger, independently selected panels and experimentally anchored target/network data \\
Parent-study MD & Ligand retention under the declared geometric criterion and system-specific endpoint assessment & Replicated, membrane-aware simulations and orthogonal binding or transport assays \\
Set-C MD pilot & Secondary trajectory-derived structural stress test for PP-01 and PP-02; single replicate per system and non-equivalent to docking-RRS & Independent mutant/wild-type replicates with declared convergence and uncertainty analyses \\
Biological polypharmacology & Not measured by the present computational scores & Target engagement plus parasite-phenotype and mechanism-of-action experiments \\
\bottomrule
\end{tabularx}
\end{table}
